% KẾT LUẬN - quy định mục 2: nêu kết quả chính và đánh giá kết quả đạt được
\chuongkhongso{Kết luận}

Báo cáo đã trình bày, tái hiện và phân tích phương pháp Future-Guided Learning (FGL) - một
khung chưng cất tri thức giữa một mô hình teacher nhìn vào dữ liệu gần tương lai hơn và một mô
hình student giải bài toán khó thực sự cần dùng - trên hai bài toán: dự báo hồi quy chuỗi
hỗn loạn Mackey-Glass và dự báo cơn động kinh từ điện não đồ CHB-MIT.

Trên Mackey-Glass, kết quả chính đạt được gồm: tái hiện thành công thực nghiệm sau khi phát
hiện và sửa tám lỗi cụ thể trong mã nguồn công khai của tác giả (từ các tệp còn nguyên dấu xung
đột merge khiến mã không chạy được, đến một khối tham số dòng lệnh bị vô hiệu do nằm trong
chuỗi chú thích); kết quả tái hiện trên sáu seed cho mức giảm MSE trung bình 26,42\%, khớp
hướng và độ lớn với mức 23,4\% bài báo công bố, ổn định trên toàn bộ các seed thử nghiệm. Trên
CHB-MIT, báo cáo đã phân tích chi tiết cơ chế FGL hiện thực hóa ``tương lai'' bằng cách định
nghĩa nhãn khác nhau cho teacher và student, và tái hiện đầy đủ cả ba giai đoạn trên ba bệnh
nhân (1, 3, 5) bằng cách chạy lại trên Kaggle. Mức tăng AUC-ROC (Baseline → Student) trung
bình $+31,2\%$ - thấp hơn mức 44,8\% bài báo công bố - và khác nhau nhiều giữa bệnh nhân:
gần như không cải thiện ở bệnh nhân Baseline đã mạnh (bệnh nhân 1), nhưng cải thiện mạnh ở hai
bệnh nhân Baseline yếu (3, 5). Điểm tích cực là chỉ số AUC bất thường (0,28) của Teacher ở lần
chạy trên máy cá nhân không lặp lại trên Kaggle.

Về hạn chế, báo cáo nhìn nhận rõ ba điểm. Thứ nhất, mức cải thiện AUC-ROC trên CHB-MIT không
ủng hộ rõ tuyên bố của bài báo và không đồng đều giữa bệnh nhân, có thể do mã nguồn CHB-MIT
không cố định seed ngẫu nhiên (khiến kết quả dao động mạnh giữa các lần chạy) và do báo cáo
mới tái hiện được ba bệnh nhân với số trial còn nhỏ (đặc biệt hai bệnh nhân 3, 5 mới có một
lần chạy). Thứ hai, các kết quả trên Mackey-Glass mỗi seed chỉ chạy một lần, nên chưa tách
được phương sai do khởi tạo ngẫu nhiên khỏi phương sai do seed chia dữ liệu, cần xem là ``khớp
hướng và độ lớn'' với bài báo hơn là một khẳng định thống kê chặt. Thứ ba, bộ dữ liệu AES hoàn
toàn chưa được tái hiện, chỉ được phân tích về lý thuyết và mã nguồn.

Hướng phát triển tiếp theo gồm: chạy lại nhiều lần hơn cho bệnh nhân 3 và 5 (hiện mới có một
lần chạy mỗi bệnh nhân) và mở rộng sang nhiều bệnh nhân khác để có ước lượng ổn định hơn về mức
cải thiện AUC-ROC, kiểm chứng trực tiếp giả thuyết về nguyên nhân chỉ số AUC bất thường ở lần
chạy trên máy cá nhân (mục~\ref{item:shuffle-false}); nếu còn thời gian, tải và tái hiện bộ dữ
liệu AES để hoàn thiện cả ba bài toán mà bài báo gốc đề cập. Báo cáo cuối kỳ sẽ bổ sung phần
cải tiến tự đề xuất (so sánh kiến trúc RNN/LSTM/GRU trên Mackey-Glass, LSTM hai chiều cho
CHB-MIT) và phần liên hệ thực tế, hiện chưa đưa vào báo cáo giữa kỳ này.
